\subsection{Segurança com Ascon}
A camada de segurança foi desenvolvida em cinco atividades, cujos resultados estão na Seção~\ref{sec:resultados_ascon}. A biblioteca \textit{ascon-suite} foi integrada ao ambiente do ESP32 e, sobre ela, foi escrita a biblioteca \textit{AsconSecure}, que cifra e autentica o \textit{payload} de cada mensagem antes de entregá-lo ao rádio. A implementação foi validada pelos \textit{Known Answer Tests} oficiais, de forma exaustiva no computador e com um subconjunto na própria placa. A chave de cada nó passou a ser gerada no próprio dispositivo durante o provisionamento, pelo gerador de \textit{hardware} do ESP32 com o CTR\_DRBG da \textit{mbedTLS}, seguindo as recomendações do NIST. Uma bateria de adulteração com 4.907 casos, sobre o texto cifrado, os dados associados, a \textit{tag} e o \textit{nonce}, confirmou que toda mensagem alterada é rejeitada. Por fim, mediram-se o tempo de execução em ciclos de \textit{clock}, a memória de \textit{heap}, pilha e \textit{Flash}, o tempo constante e o comparativo com o AES-128-GCM nas duas placas, e o consumo de energia com resistor \textit{shunt} e osciloscópio, também nas duas placas.

%%%%%%%%%%%%%%%%%%%%%%%%%%%%%%%%%%%%%%%%%%%%%%%%%%%%
